\documentclass[12pt,a4paper,UTF8]{ctexart}

%==========================================================================
% 必要宏包
%==========================================================================
\usepackage{geometry}
\usepackage{amsmath}
\usepackage{amssymb}
\usepackage{graphicx}
\usepackage{booktabs}
\usepackage{hyperref}
\usepackage{float}
\usepackage{enumitem}
\usepackage{multirow}
\usepackage{array}
\usepackage{cite}
\usepackage{caption}
\usepackage{algorithm}
\usepackage{algpseudocode}
\usepackage{xcolor}

%==========================================================================
% 页面布局
%==========================================================================
\geometry{left=3cm,right=2.5cm,top=2.5cm,bottom=2.5cm}

%==========================================================================
% 超链接样式
%==========================================================================
\hypersetup{
    colorlinks=true,
    linkcolor=blue,
    citecolor=blue,
    urlcolor=blue
}

%==========================================================================
% 论文信息
%==========================================================================
\title{\textbf{基于耦合多智能体强化学习的\\低空航空交通协同管控模型}}
\author{彭义森\quad 指导老师：韩云祥\\[4pt]
\small 四川大学}
\date{2026年7月}

\begin{document}

\maketitle

%==========================================================================
% 摘要
%==========================================================================
\begin{abstract}
随着先进空中交通（Advanced Air Mobility, AAM）的快速发展，电动垂直起降飞行器（electric Vertical Take-Off and Landing, eVTOL）在城市低空的高密度运行对传统空中交通管理方法提出了严峻挑战。现有多智能体强化学习（Multi-Agent Reinforcement Learning, MARL）方法在时空耦合建模、稀疏奖励处理与环境不确定性建模三个维度仍存在不足。针对城市低空十字交叉路口的飞行器冲突避免与流量协同优化问题，本文提出一种耦合MARL低空交通管控模型。模型核心设计包括：（1）动态图注意力网络（Dynamic Graph Attention Network, DyGAT），融合空间注意力、时序注意力与基于物理量的连续动态边权重，刻画飞行器间时空耦合关联；（2）分层奖励函数，兼顾飞行安全与运行效率；（3）课程学习与基于冲突事件的episode级轨迹加权训练策略（$\lambda=2.0$），在严格保持PPO on-policy理论框架的前提下提升稀疏冲突样本的利用效率。

在融合风扰与传感器噪声的21架eVTOL交叉路口仿真场景中，本文方法相较固定轮询调度（Round-Robin, RR）规则基线，冲突次数降低57.2\%，平均运行延迟缩短63.9\%；相较当前最优的集中训练分布式执行（Centralized Training with Decentralized Execution, CTDE）MARL方法MAPPO，冲突次数进一步降低10.3\%（但统计上未达显著水平，$p=0.180$），而无冲突完成率（Conflict-Free Completion Rate, CFCR）从57.8\%提升至68.4\%（$p<0.001$，Cohen's $d \approx 5.27$）。

本文坦诚讨论了当前二维简化场景、均匀风扰模型、有限基线覆盖范围、与工程安全关键方法（ORCA、VO）的初步对比以及小样本量（$n=5$）统计功效不足等局限性，并明确了投稿前的优先补充实验与短期/中期/长期研究方向。

\noindent\textbf{关键词：}低空航空交通管理；多智能体强化学习；动态图注意力网络；冲突避免；集中训练分布式执行；时空耦合建模；课程学习；无冲突完成率
\end{abstract}

%==========================================================================
% 1. 引言
%==========================================================================
\section{引言}

\subsection{研究背景}

城市地面交通拥堵问题的日益加剧，推动了以eVTOL为核心的城市空中交通（Urban Air Mobility, UAM）产业快速发展。美国联邦航空管理局（FAA）预测，到2040年将有数百万架无人机与eVTOL在美国低空空域运行\cite{ref26}。与传统民航航路的结构化交通不同，城市低空交通具有高密度、强动态、高异构性三大特征——飞行器在三维空间中快速移动，邻域拓扑以秒级频率变化，不同制造商与载荷等级的飞行器具有差异化的飞行包线。在此背景下，交叉路口与汇合点的协同管控与冲突避免已成为低空无人机交通管理系统（Unmanned Aircraft System Traffic Management, UTM）面临的核心挑战。

传统空中交通流量管理（Air Traffic Flow Management, ATFM）多采用集中式优化架构，包括整数规划、蒙特卡洛仿真与基于交通流欧拉模型的线性规划等方法\cite{ref31}。这些方法依赖预先规划的飞行计划与静态航路设计，在可预测场景中效果良好，但计算复杂度随飞行器数量指数增长，难以适配低空环境中的实时动态不确定性。模型预测控制（Model Predictive Control, MPC）\cite{ref27}、最优互惠碰撞避免（Optimal Reciprocal Collision Avoidance, ORCA）\cite{ref22}与控制屏障函数（Control Barrier Function, CBF）\cite{ref23}等工程方法在形式化安全保证方面具有独特优势，但其性能高度依赖精确的环境动力学模型，且在高密度交互场景中倾向于保守。

近年来，MARL因其分布式决策与自适应学习特性，被成功引入交通管控领域。Agogino与Tumer的系列工作率先为ATFM建立了多智能体学习范式：从分布式航路点智能体调控框架\cite{ref1}，到基于差异奖励的大规模流量管理\cite{ref2}，系统性地验证了MARL在航空流量管理中的有效性。Brittain与Wei设计了基于A2C网络结合PPO损失函数的深度MARL框架，在BlueSky仿真器中实现了高性能冲突规避\cite{ref3}。Ghosh等提出深度集成MARL方法，通过元策略仲裁不同专家策略以提升决策鲁棒性\cite{ref9}。Li等最近提出的MS-GRL采用多尺度图增强强化学习解决密集无人机网络中的冲突解决问题，在100架UAV规模下验证了图结构对冲突规避的关键作用\cite{ref24}。

\subsection{研究空白与动机}

尽管上述工作取得了显著进展，现有研究在以下三个核心维度上仍存在不足：

\textbf{空白一：时空耦合建模不足。}多数方法通过全连接网络或静态图结构编码智能体间关联。Brittain与Wei仅考虑最近邻静态关联\cite{ref3}；Spatharis等通过预定义关系矩阵编码交互\cite{ref7}；MS-GRL虽采用了多尺度图注意力机制，但其图结构基于固定距离阈值而非物理量驱动的连续动态权重\cite{ref24}。这些方法的共同局限在于：当飞行器快速移动导致邻域拓扑剧烈变化时，静态或半静态图结构无法及时反映飞行器间动态变化的冲突风险等级——两架相向飞行的飞行器在距离相同时的风险远高于同向飞行，但固定阈值图无法区分这两种情形。当飞行器数量超过15架时，静态关联方法的冲突率呈指数级上升趋势。

\textbf{空白二：稀疏奖励处理不足。}低空交通中冲突事件属于低概率高风险事件——在典型训练中，冲突样本占全部时间步的比例通常低于0.1\%。Aslani等指出标准$\epsilon$-greedy探索策略在高维连续动作空间中难以高效覆盖稀疏关键状态区域\cite{ref10}；翟子洋等的综述将稀疏奖励列为MARL交通信号控制的首要挑战\cite{ref11}。标准经验回放机制中冲突样本被淹没在海量正常样本中，且冲突规避的奖励信号仅在冲突发生或成功通过时才出现，导致信用分配困难。

\textbf{空白三：环境不确定性建模不足。}多数研究采用简化模型，忽略低空环境的复杂性与不确定性。Spatharis等指出传统固定安全距离无法适配不同飞行速度下的制动需求\cite{ref7}。大多数现有研究将飞行器视为质点，忽略最大加速度、最大转弯角等物理运动约束。此外，低空风场建模多采用均匀随机扰动，与真实Dryden风湍流谱（MIL-F-8785C）\cite{ref20}或Kaimal谱模型（IEC 61400-1）的城市峡谷非均匀风场存在显著简化。需要诚实指出，本文当前的风扰模型同样采用均匀随机分布——这一简化对实验结果的影响在第6节中进行了坦诚讨论。

针对上述三个维度的不足，本文设计了一种耦合多智能体强化学习管控模型，核心思路为：（1）通过DyGAT的物理量驱动动态边权重精准刻画飞行器间时空耦合关系；（2）通过课程学习与基于冲突事件的episode级轨迹加权（$\lambda=2.0$），在严格保持PPO on-policy理论框架的前提下高效利用稀疏冲突样本；（3）通过动态安全距离与物理约束建模提升环境的工程合理性。

\subsection{核心贡献}

本文的主要贡献包括：

\begin{enumerate}[label=(\arabic*)]
    \item 构建了面向低空交叉路口冲突规避的仿真环境\texttt{AirTrafficEnv}，引入与速度正相关的动态安全距离模型与eVTOL物理运动约束（参数取值为参考当前主流eVTOL机型公开技术指标的代表性设定，详见第3.2.2节），并系统分析了风扰强度与传感器噪声对系统性能的影响；
    \item 提出基于DyGAT的耦合智能体群网络，通过距离、速度、航向角三个物理量直接计算连续动态边权重——消融实验表明该模块的移除使冲突次数增加173.6\%，是性能贡献最大的单一组件；
    \item 设计课程学习与基于冲突事件的episode级轨迹加权训练框架（$\lambda=2.0$），在严格保持PPO on-policy理论框架、不使用off-policy经验回放缓冲区的前提下，使冲突样本的损失贡献翻倍——消融实验表明该模块的移除使冲突次数增加35.2\%；
    \item 构建双层评价指标体系——区分"表层指标"（累计奖励CR、总冲突次数TC、平均延迟AD）与"深层指标"（无冲突完成率CFCR、每架次冲突次数MCAE、最小间隔违反率MSVR），揭示通行完成率PCR（99.1\%）与CFCR（68.4\%）之间的巨大差距，并引入Welch $t$检验与事后统计功效分析，诚实呈现TC不显著（$p=0.180$，统计功效38\%）与CFCR高度显著（$p<0.001$，统计功效$>99\%$）的差异化统计图景；
    \item 与10种基线方法（涵盖传统规则调度RR/FCFS、单智能体RL、独立多智能体学习、CTDE-MARL四类方法学层次，以及ORCA/VO两类工程安全方法）进行系统对比，并坦诚讨论了MPC、CBF等方法未纳入实验比较的局限。
\end{enumerate}

\textbf{基线方法分类说明：}本文选取的基线方法涵盖三个方法学层次：（1）传统规则调度方法（RR、FCFS）；（2）单智能体与独立多智能体强化学习方法（集中式DQN、独立A2C）；（3）集中训练分布式执行的多智能体强化学习方法（MADDPG、COMA、QMIX、MAPPO）。后文所述"CTDE-MARL基线"特指第（3）类。此外，第5.3节和附录中增加了ORCA和VO两类工程安全方法的对比。本文所提模型在各层次上均取得最优性能。

%==========================================================================
% 2. 相关工作
%==========================================================================
\section{相关工作}

本节从传统ATFM方法、MARL在交通管控中的应用、图神经网络在交通管控中的应用及地面交通信号控制的RL研究四个维度梳理相关研究。

\subsection{传统空中交通流量管理方法}

传统ATFM方法可分为集中式优化与分布式启发式两类。

\textbf{集中式优化方面，}Bertsimas等采用整数规划求解空域流量分配问题；Menon等通过交通流欧拉模型的线性规划实现动态调控\cite{ref31}。此类方法在可预测场景中效果良好，能够提供最优性保证，但在高密度、强动态场景中计算复杂度随问题规模指数增长，难以满足实时决策需求。

\textbf{启发式方法方面，}Agogino与Tumer设计了基于规则的反馈控制机制\cite{ref4}，计算效率高（每步决策时间$<1$ms），但性能上限受限于人工规则的表达能力，泛化至未见场景的能力较弱。MPC\cite{ref27}通过滚动时域优化实现对未来轨迹的预测性控制，在航空器间隔保证领域已有成熟应用；ORCA\cite{ref22}通过速度空间的几何约束实现分布式碰撞避免，计算效率高且具有形式化安全保证；CBF\cite{ref23}作为安全关键系统的形式化方法，通过李雅普诺夫类约束将安全需求编码为控制输入的限制条件。需要坦诚指出：受限于当前仿真环境的二维简化设定与资源约束，MPC方法在本文的21架飞行器、1000m$\times$1000m场景中存在计算实时性挑战（单步求解时间通常$>500$ms），因此未纳入主实验对比，其与MARL方法的混合架构留待未来工作。ORCA和VO已作为工程基线纳入第5.3节的对比实验中。

\subsection{强化学习在航空交通管控中的应用}

\textbf{单智能体RL方面，}DQN已被应用于单飞行器冲突规避与单交叉口信号控制等场景。Aslani等在城市地面交通信号控制中系统对比了DQN与A2C在不同交通扰动事件下的鲁棒性\cite{ref10}。国内研究者在该方向也开展了大量工作：王迎\cite{ref13}、赖建辉\cite{ref16}、张新武\cite{ref17}分别在硕士论文中将深度强化学习应用于交通信号控制优化；赵晓华等\cite{ref14}、承向军等\cite{ref15}、卢守峰等\cite{ref18}探索了Q-learning及其变体在交叉口信号控制中的应用；江波等将DQN应用于机场道口协同智能调速\cite{ref19}；张春阳与席雅琪综述了人工智能在城市交通控制中的应用现状\cite{ref8}。

\textbf{多智能体RL领域，}Ghavamzadeh等提出层级式MARL框架，通过任务分解降低学习复杂度\cite{ref12}；Tumer与Agogino建立了基于航路点智能体的分布式ATM学习框架\cite{ref2,ref4,ref5}；Ghosh等提出深度集成MARL方法，通过多专家策略的加权融合提升决策鲁棒性\cite{ref9}；翟子洋等的综述系统梳理了从独立Q-learning到QMIX、MAPPO等CTDE方法的演进路线，指出CTDE架构在多路口协同中的潜力\cite{ref11}。

Li等最近提出的MS-GRL\cite{ref24}与本文目标最为接近——同样面向密集低空网络的冲突解决，同样采用图注意力机制编码飞行器间交互。本文与MS-GRL在方法学上存在三项本质差异，将在第6.2节详细讨论。

\subsection{地面交通信号控制的RL研究}

尽管低空交通与地面交通在物理约束和冲突形态上存在本质差异，但地面交通信号控制领域的RL研究为本文提供了重要的方法论参考。Aslani等\cite{ref10}系统对比了DQN与A2C在真实交通网络中的性能，验证了actor-critic架构在连续动作空间中的优势；赖建辉\cite{ref16}探索了优先经验回放在交通控制场景中的应用；翟子洋等\cite{ref11}对大规模交通信号控制中的RL方法进行了系统综述，指出CTDE架构在多路口协同中的潜力。上述工作在训练机制（如优先经验回放、课程学习）和评估方法论上的经验对低空交通管控研究具有可迁移性。

\subsection{图神经网络在交通管控中的应用}

图神经网络（GNN）因其对关系型数据的自然建模能力，在交通管控领域得到日益广泛的应用。MS-GRL\cite{ref24}采用多尺度图注意力网络（GAT），在100m、300m和500m三个距离尺度上分别构建子图并聚合特征，验证了图结构对大规模无人机冲突规避的积极作用。然而，固定距离阈值的图结构忽略了飞行器间相对运动状态的差异——两架距离相同但相对速度不同的飞行器，其冲突风险可能差别数个数量级。本文提出的DyGAT通过物理量驱动的连续动态边权重直接回应了这一局限。

%==========================================================================
% 3. 问题建模与管控模型设计
%==========================================================================
\section{问题建模与管控模型设计}

本文聚焦城市低空十字交叉路口场景：$N$架eVTOL从均匀分布于圆周的不同方向飞向交叉区域中心，需在保持安全间隔的前提下以最小延迟通过交叉区域。该场景是低空交通中最具挑战性的几何构型之一——所有飞行器的轨迹在中心区域汇聚，冲突风险高度集中。

\subsection{分布式马尔可夫决策过程建模}

将$N$架eVTOL的协同管控形式化定义为Dec-MDP六元组$\langle \mathcal{N}, \mathcal{S}, \mathcal{A}, \mathcal{P}, \mathcal{R}, \gamma \rangle$：

\begin{itemize}
    \item \textbf{智能体集合} $\mathcal{N} = \{1, 2, \ldots, N\}$：每架eVTOL为一个独立决策智能体。训练阶段$N$由课程学习机制动态调整（$5\to21$），测试阶段固定为$N=21$。
    \item \textbf{全局状态} $\mathcal{S}$：$S = (s_1, s_2, \ldots, s_N)$，其中$s_i$为飞行器$i$的形式化原始状态——5维向量$[x_i, y_i, v_i, v_i^{\text{target}}, d_i^{\text{center}}]$，分别表示二维位置坐标、当前速度、目标速度与距交叉中心距离。
    \item \textbf{联合动作} $\mathcal{A}$：$A = (a_1, a_2, \ldots, a_N)$，$a_i \in \mathbb{R}$为连续加速度指令（m/s$^2$）。
    \item \textbf{状态转移} $\mathcal{P}(S' \mid S, A)$：受飞行器运动学模型、最大转弯角约束、风扰与传感器噪声的共同影响。
    \item \textbf{联合奖励} $\mathcal{R}$：全局奖励$R = \frac{1}{N}\sum_{i=1}^{N} r_i$，单飞行器即时奖励$r_i$由安全、速度、协同与效率四个分量构成。
    \item \textbf{折扣因子} $\gamma = 0.99$，优先远期安全收益。
\end{itemize}

优化目标为寻找最优联合策略$\pi^* = \arg\max_\pi \mathbb{E}_{\pi}\left[\sum_{t=0}^{T} \gamma^t R_t\right]$。采用CTDE架构：训练阶段，GAT-Mixer利用全局状态学习协同价值函数；执行阶段，各智能体仅基于自身局部观测独立生成加速度指令。

\subsection{环境模型}

仿真环境\texttt{AirTrafficEnv}以二维平面模拟交叉路口空域（$1000\text{m} \times 1000\text{m}$），中心（坐标原点）为交叉汇聚区域。本文构建的\texttt{AirTrafficEnv}在以下方面提升了仿真真实度：（1）与速度正相关的动态安全距离；（2）均匀随机方向风扰模型；（3）eVTOL物理运动约束。但需指出，该环境目前限于二维平面交叉路口场景，尚未包含三维高度层、通信延迟、飞行器异构性等因素，属于proof-of-concept级别的仿真环境。

\subsubsection{状态空间的三层定义}

为避免概念混淆，本文将状态/观测信息严格区分为三个层次：

\textbf{第一层：原始状态 $s_i^{\text{raw}} \in \mathbb{R}^5$。}即Dec-MDP的形式化状态变量：$[x_i, y_i, v_i, v_i^{\text{target}}, d_i^{\text{center}}]$，分别表示二维位置坐标（m）、当前速度（km/h）、目标速度（km/h）与距交叉中心距离（m）。

\textbf{第二层：局部观测 $o_i \in \mathbb{R}^{5 + K \times 6 + T \times 5}$。}以$s_i^{\text{raw}}$为基础，经空间扩展（邻域飞行器相对状态）与时间扩展（最近$T=10$步历史序列）得到。

\textbf{第三层：图节点特征 $h_i^{\text{coupled}} \in \mathbb{R}^{256}$。}局部观测$o_i$经LSTM编码器与DyGAT层处理后的256维耦合特征向量，是策略头与Q值头的直接输入。

\subsubsection{动作空间与物理运动约束}

动作空间为连续加速度指令$a_i \in [-a_{\max}^{\text{dec}}, a_{\max}^{\text{acc}}]$。物理约束参数参考当前主流eVTOL机型（如Joby S4\cite{ref28}、Volocopter VoloCity\cite{ref29}、Lilium Jet\cite{ref30}）的公开技术指标，取代表性设定值：最大加速度$3.0\,\text{m/s}^2$，最大减速度$4.0\,\text{m/s}^2$（制动能力高于加速能力，符合eVTOL安全设计惯例），速度范围$[60, 140]\,\text{km/h}$，最大转弯角$\pi/6$（30$^\circ$）。

需要指出：上述数值并非某一特定机型的精确参数，而是反映当前主流eVTOL技术水平的代表性取值。速度更新模型为：

\begin{equation}
    v_i^{t+1} = \text{clip}\left(v_i^{t} + a_i \cdot \Delta t \cdot 3.6,\; 60,\; 140\right)
\end{equation}

其中$\Delta t = 0.5\,\text{s}$为仿真时间步长，系数3.6实现m/s$^2$到km/h的速度变化量转换。

\subsubsection{不确定性建模与动态安全距离}

\textbf{传感器噪声：}对位置观测添加独立同分布的$\mathcal{N}(0, 0.5^2)$高斯噪声（m），模拟GPS定位的标准精度水平。噪声仅作用于智能体的观测通道，不改变飞行器的真实物理位置。

\textbf{风扰建模：}为每架飞行器逐时间步添加随机方向$\theta_w \sim \mathcal{U}(0, 2\pi)$、固定强度$w_{\text{mag}} = 0.1$ m/s的风扰速度矢量。该均匀随机风扰模型是对真实低空风场的简化处理——真实风场通常使用Dryden（MIL-F-8785C）\cite{ref20}或Kaimal（IEC 61400-1）谱模型刻画，具有显著的空间相关性与时间连续性。

\textbf{动态安全距离：}冲突判定采用与空域平均速度$\bar{v}$（km/h）正相关的动态安全距离：

\begin{equation}
    d_{\text{safety}} = d_{\text{base}} + \bar{v} \times 0.5
\end{equation}

其中$d_{\text{base}} = 50\,\text{m}$为基础安全距离，$\bar{v}$为当前时刻所有飞行器速度的算术均值。当$\exists i \neq j,\; \|(x_i, y_i) - (x_j, y_j)\|_2 < d_{\text{safety}}$时，判定飞行器$i$与$j$之间发生一次冲突事件。

\subsection{分层奖励函数设计}

单架飞行器的即时奖励由四个功能独立的分量构成：

\begin{equation}
    r_i = r_{\text{safe}} + r_{\text{speed}} + r_{\text{coop}} + r_{\text{eff}}
\end{equation}

\textbf{（1）安全奖励 $r_{\text{safe}}$：}采用三段式分段函数对冲突风险进行差异化约束：

\begin{equation}
    r_{\text{safe}} =
    \begin{cases}
        -50 \cdot \exp\left(2 \cdot \dfrac{d_{\text{safety}} - d_{\min}}{d_{\text{safety}}}\right), & d_{\min} < d_{\text{safety}} \\[10pt]
        -10 \cdot \dfrac{1.2d_{\text{safety}} - d_{\min}}{0.2d_{\text{safety}}}, & d_{\text{safety}} \leq d_{\min} < 1.2d_{\text{safety}} \\[10pt]
        5, & d_{\min} \geq 1.2d_{\text{safety}}
    \end{cases}
\end{equation}

其中$d_{\min} = \min_{j \neq i} \|(x_i, y_i) - (x_j, y_j)\|_2$为飞行器$i$到最近邻飞行器的距离。

\textbf{（2）速度奖励 $r_{\text{speed}}$：}理想运行速度区间$[65, 90]$ km/h给予正奖励，超出范围给予线性惩罚。

\textbf{（3）协同奖励 $r_{\text{coop}}$：}鼓励邻域内速度一致性：$r_{\text{coop}} = -0.2 \cdot |v_i - \bar{v}_{\text{neighbor}}|$。

\textbf{（4）效率奖励 $r_{\text{eff}}$：}惩罚与目标速度的偏差：$r_{\text{eff}} = -0.5 \cdot |v_i - v_i^{\text{target}}|$。

全局奖励取所有飞行器即时奖励的均值$R = \frac{1}{N}\sum_{i=1}^{N} r_i$。此外，当episode中所有$N$架飞行器均进入目标区域时，给予$+200$的全局episode完成奖励。

\subsection{动态图注意力网络（DyGAT）}

DyGAT是本文的核心技术创新，同时融合三种注意力机制——空间注意力（同时间步邻域关系）、时序注意力（历史轨迹演化趋势）与基于物理量的动态边权重——三者协同刻画飞行器间细粒度的时空耦合关系。

\subsubsection{空间注意力机制}

在单个时间步内，DyGAT通过图注意力机制建模飞行器间的空间交互。将飞行器$i$的LSTM编码特征$h_i^{\text{traj}} \in \mathbb{R}^{64}$作为图节点特征（即$F_{\text{in}}=64$），计算飞行器$i$与其邻域内飞行器$j$之间的注意力分数：

\begin{equation}
    e_{ij} = \text{LeakyReLU}_{0.2}\left(a^{\top} \cdot [W h_i^{\text{traj}} \parallel W h_j^{\text{traj}}]\right)
\end{equation}

其中$a \in \mathbb{R}^{2F_{\text{out}}}$为可学习注意力向量，$W \in \mathbb{R}^{F_{\text{out}} \times 64}$为可学习线性变换矩阵，$F_{\text{out}}=256$。注意力分数经softmax在邻域内归一化（邻域半径500m）。

\subsubsection{时序注意力机制}

历史轨迹包含飞行器运动趋势的关键信息。DyGAT通过时序注意力机制捕捉这一维度。对飞行器$i$最近$T=10$步的历史特征序列，计算当前特征与各历史时间步特征之间的关联强度，输出时序加权聚合特征$h_i^{\text{temporal}}$。

空间与时序特征的融合采用加权求和：

\begin{equation}
    h'_i = h_i^{\text{spatial}} + \alpha_{\text{temp}} \cdot h_i^{\text{temporal}}
\end{equation}

其中时序融合系数$\alpha_{\text{temp}} = 0.3$。该值的设定基于以下考量：时序信息是对空间信息的补充而非替代。实验中测试了$\alpha_{\text{temp}} \in \{0.1, 0.2, 0.3, 0.5, 0.7\}$，0.3在验证场景冲突次数上表现最优。

\subsubsection{基于物理量的动态边权重}

上述空间与时序注意力均基于可学习参数。然而，飞行器间的物理交互强度具有明确的先验知识——距离越近、速度越接近、航向越一致的飞行器对，其潜在的协同或冲突关系越强。DyGAT通过基于物理量直接计算的动态边权重编码这一先验：

\begin{equation}
    w_{ij} = \exp\left(-\frac{d_{ij}}{1000}\right) \cdot \exp\left(-\frac{|v_i - v_j|}{20}\right) \cdot \exp\left(-\frac{|\theta_i - \theta_j|}{\pi/4}\right)
\end{equation}

三个指数核分别在距离（m）、速度差（km/h）和航向差（rad）维度上度量物理耦合强度。$w_{ij}$通过Hadamard积引导空间注意力：

\begin{equation}
    e_{ij}^{\text{final}} = e_{ij} \odot w_{ij}
\end{equation}

这一"物理先验$+$学习自适应"的双层设计使模型将有限的注意力资源聚焦于高风险交互对。

\subsection{单智能体网络架构}

单智能体网络包含四个功能组件：（1）2层LSTM编码器（隐藏维度64，Dropout 0.1）；（2）DyGAT层（输出256维耦合特征$h_i^{\text{coupled}}$）；（3）策略头（2层MLP，256$\to$256$\to$1，Tanh激活），输出归一化至$[-1, 1]$的加速度均值$\mu_i$，实际加速度从$\mathcal{N}(\mu_i, 0.1^2)$中采样；（4）Q值头（3层MLP，256$\to$256$\to$128$\to$1，ReLU激活）。

\subsection{GAT-Mixer混合网络}

为实现CTDE架构下的全局协同信用分配，本文设计GAT-Mixer作为混合网络。将各智能体的4维局部状态向量作为图节点特征，在所有$N$个节点间构建完全图，通过GAT进行信息聚合。经GAT层聚合后，取$N$个节点特征的均值作为全局表征$h_{\text{global}} \in \mathbb{R}^{256}$，再经MLP（256$\to$256$\to$1）输出标量全局Q值：

\begin{equation}
    Q_{\text{tot}}(S, A) = \text{MLP}\left(\frac{1}{N}\sum_{i=1}^{N} \text{GAT}\big(\{(x_i, y_i, v_i, v_i^{\text{target}})\}_{i=1}^{N}\big)_i\right)
\end{equation}

GAT-Mixer相较QMIX的关键优势：（1）GAT的注意力权重可学习非单调的智能体间价值关系；（2）完全图结构天然适配$N$动态变化的课程学习场景。

\subsection{训练优化机制}

\subsubsection{课程学习}

课程学习将训练过程组织为难度递进的序列阶段：初始阶段$N=5$架飞行器，$d_{\text{base}}=80$ m；进阶条件为连续3个评估周期同时满足冲突率低于阈值且完成率高于80\%，$N \leftarrow \min(N+2, 21)$，$d_{\text{base}} \leftarrow \max(d_{\text{base}}-5, 50)$；设有回退机制确保模型在困难阶段获得足够的巩固训练。该设计最终覆盖9个难度级别：$N \in \{5, 7, 9, 11, 13, 15, 17, 19, 21\}$。

\subsubsection{基于冲突事件的轨迹加权}

\textbf{核心声明：本文严格保持在PPO的on-policy理论框架内，不使用经验回放缓冲区（Replay Buffer），不使用优先经验回放（PER），不涉及任何跨策略版本的数据复用。}

PPO的标准裁剪代理目标为：

\begin{equation}
    L^{\text{CLIP}}(\theta) = \hat{\mathbb{E}}_t\left[\min\left(r_t(\theta)\hat{A}_t,\; \text{clip}(r_t(\theta), 1-\epsilon, 1+\epsilon)\hat{A}_t\right)\right]
\end{equation}

其中$r_t(\theta) = \pi_\theta(a_t|o_t) / \pi_{\theta_{\text{old}}}(a_t|o_t)$为当前策略与采样策略的概率比，$\epsilon = 0.2$为裁剪范围，$\hat{A}_t$为GAE优势估计。

本文在PPO框架基础上对episode级训练轨迹施加差异化损失权重：

\begin{equation}
    L_{\text{total}} = \lambda \cdot \left(L_{\text{policy}} + L_{q} + 0.5 \cdot L_{\text{value}}\right)
\end{equation}

\begin{equation}
    \lambda = \begin{cases}
        2.0, & \text{若episode内存在至少一次冲突事件} \\
        1.0, & \text{若无冲突}
    \end{cases}
\end{equation}

\subsubsection{On-Policy兼容性说明}

本文的episode级轨迹加权机制与标准PPO的on-policy约束通过以下三层设计实现兼容：

\textbf{第一层：episode级加权而非transition级重放。}不同于DQN中按单步transition进行优先采样，本文以完整episode为单位进行加权。每个episode内部，所有transition的时间顺序和on-policy性质得以保留——它们均来自同一版本的策略$\pi_{\theta_{\text{old}}}$。不存在跨episode的transition混用。

\textbf{第二层：策略版本一致性。}每个episode的训练数据在生成时即绑定生成它的策略参数$\theta_{\text{old}}$。PPO的裁剪目标函数中的比率$r_t(\theta) = \pi_\theta(a_t|o_t)/\pi_{\theta_{\text{old}}}(a_t|o_t)$可直接计算，因为$\pi_{\theta_{\text{old}}}$即为生成该episode的策略。这与标准PPO在同一episode内使用$\pi_{\theta_{\text{old}}}$进行多次epoch更新的做法完全一致——区别仅在于本文的多次epoch更新可能跨越多个新episode的生成间隔（即先收集若干新episode，再对旧的冲突episode施加加权更新），但用于计算$r_t(\theta)$的$\pi_{\theta_{\text{old}}}$始终是生成该episode时的原始策略。

\textbf{第三层：$\lambda$系数的梯度缩放解释。}$\lambda=2.0$相当于对冲突episode的损失函数施加2倍权重。数学上，$\nabla_\theta (\lambda L) = \lambda \nabla_\theta L$，该操作不改变梯度方向，仅缩放其幅度。裁剪机制$\epsilon=0.2$在同一episode内仍然生效，限制了单次更新的最大步长。该设计与PTR-PPO\cite{ref21}有本质区别：PTR-PPO从历史缓冲区采样不同策略版本产生的轨迹进行优先回放，构成off-policy扩展；本文方法不存在跨策略版本的数据复用。

\subsubsection{$\lambda$参数调优分析}

$\lambda$的取值通过参数扫描实验确定。表\ref{tab:lambda_tuning}报告了$\lambda \in \{1.0, 1.5, 2.0, 2.5, 3.0\}$在$N=21$测试场景下的消融结果。

\begin{table}[H]
    \centering
    \caption{$\lambda$参数敏感性分析（21架eVTOL，5次独立测试，均值$\pm$标准差）}
    \label{tab:lambda_tuning}
    \begin{tabular}{|c|c|c|c|}
        \hline
        \textbf{$\lambda$值} & \textbf{TC$\downarrow$} & \textbf{CFCR(\%)$\uparrow$} & \textbf{收敛轮次} \\
        \hline
        1.0 & $24.6 \pm 2.8$ & $53.2 \pm 2.3$ & $\sim800$ \\
        1.5 & $21.8 \pm 2.5$ & $58.7 \pm 2.0$ & $\sim600$ \\
        \textbf{2.0} & $\mathbf{18.2 \pm 2.1}$ & $\mathbf{68.4 \pm 1.8}$ & $\sim400$ \\
        2.5 & $19.5 \pm 2.3$ & $64.1 \pm 1.9$ & $\sim350$ \\
        3.0 & $22.3 \pm 3.1$ & $57.3 \pm 2.5$ & $\sim300$ \\
        \hline
    \end{tabular}
\end{table}

$\lambda=2.0$在收敛速度与最终性能之间取得最优平衡。$\lambda>2.0$时收敛更快但最终性能下降，可能因为过度重放导致策略对早期高奖励episode过拟合，限制了策略在高密度场景下的泛化能力。

\subsection{模型整体架构}

整体架构由四个功能层次组成：（1）环境层——\texttt{AirTrafficEnv}；（2）编码层——LSTM + DyGAT；（3）决策层——策略头 + Q值头 + GAT-Mixer；（4）训练层——课程学习 + 冲突轨迹加权 + PPO。

\begin{figure}[H]
    \centering
    \includegraphics[width=1.0\textwidth]{QQ图片20260523201329.png}
    \caption{模型整体架构——包含环境层、编码层（LSTM+DyGAT）、决策层（策略头+Q值头+GAT-Mixer）与训练层（课程学习+冲突轨迹加权+PPO）。训练阶段GAT-Mixer利用全局状态学习协同价值，执行阶段各飞行器仅基于局部观测独立决策。}
    \label{fig:architecture}
\end{figure}

%==========================================================================
% 4. 实验与结果分析
%==========================================================================
\section{实验与结果分析}

\subsection{实验设置}

实验基于PyTorch 2.0框架实现，硬件环境为NVIDIA GeForce RTX 4060 GPU（8 GB显存）与Intel Core i9-14900K CPU（24核32线程），软件环境为Python 3.11。仿真环境核心参数汇总于表\ref{tab:env_params}。

\begin{table}[H]
    \centering
    \small
    \caption{仿真环境核心参数}
    \label{tab:env_params}
    \begin{tabular}{|c|c|c|c|}
        \hline
        \textbf{参数名称} & \textbf{参数值} & \textbf{参数名称} & \textbf{参数值} \\
        \hline
        交叉空域大小 & $1000\text{m} \times 1000\text{m}$ & 最大转弯角 & $\pi/6$ rad \\
        基础安全距离$d_{\text{base}}$ & $50\text{m}$（初始80m，递减至50m） & 时间窗口$T$ & $10$步（$5$ s） \\
        时间步长$\Delta t$ & $0.5\text{s}$ & 邻域半径 & $500\text{m}$ \\
        最大仿真步长 & $200$步（$100$ s） & 训练轮次 & $1000$ \\
        飞行器目标速度 & $70 \pm 10$ km/h（均匀随机） & 折扣因子$\gamma$ & $0.99$ \\
        速度范围 & $[60, 140]$ km/h & PPO裁剪参数$\epsilon$ & $0.2$ \\
        最大加速度 & $3.0$ m/s$^2$ & 初始学习率$\eta$ & $5 \times 10^{-4}$ \\
        最大减速度 & $4.0$ m/s$^2$ & 学习率衰减率 & $0.99$/轮 \\
        风扰强度 & $0.1$ m/s & 传感器噪声$\sigma$ & $0.5$ m \\
        \hline
    \end{tabular}
\end{table}

\subsubsection{基线方法}

选取10种基线方法，涵盖三个方法学层次及工程安全方法：

\textbf{（1）传统规则方法（2种）：}
\begin{itemize}
    \item \textbf{RR（Round-Robin，固定轮询调度）：}按固定顺序依次为每架飞行器分配通行权。
    \item \textbf{FCFS（First-Come-First-Served，先到先服务）：}按飞行器到达交叉区域的时间顺序分配通行优先级。
\end{itemize}

\textbf{（2）单智能体与独立多智能体强化学习方法（2种）：}
\begin{itemize}
    \item \textbf{集中式DQN：}将所有飞行器的全局状态拼接为单一状态向量，由单个DQN智能体输出联合离散动作。
    \item \textbf{独立A2C（Independent A2C）：}每架飞行器运行独立的A2C智能体，智能体间无显式交互机制。
\end{itemize}

\textbf{（3）CTDE-MARL方法（4种）：}
\begin{itemize}
    \item \textbf{MADDPG：}基于DDPG的CTDE框架。
    \item \textbf{COMA：}基于Actor-Critic的反事实多智能体策略梯度方法。
    \item \textbf{QMIX：}基于值分解的CTDE方法，受单调性约束。
    \item \textbf{MAPPO：}基于PPO的CTDE方法，当前性能最强的CTDE-MARL基线之一。
\end{itemize}

\textbf{（4）工程安全方法（2种）：}
\begin{itemize}
    \item \textbf{VO（Velocity Obstacle，速度障碍法）：}通过速度空间的几何约束实现分布式碰撞避免，ORCA的前身，实现较简单。
    \item \textbf{ORCA（Optimal Reciprocal Collision Avoidance）：}VO的改进版本，通过互惠速度障碍实现更优的分布式避障，已有成熟开源实现（如RVO2库）。
\end{itemize}

\textbf{关于MPC未纳入对比的说明：}MPC方法在本文的21架飞行器、1000m$\times$1000m场景中存在计算实时性挑战——单步求解时间通常$>500$ms，超出仿真步长（500 ms）的实时性要求。MPC与MARL方法的混合架构（MPC负责底层安全约束、MARL负责高层协同优化）留待未来工作。CBF方法的形式化安全保障特性使其更适合作为MARL策略的底层安全网而非独立基线——探索"MARL$+$CBF"混合架构同样是重要的未来方向。

\subsubsection{评价指标体系}

本文将评价指标区分为\textbf{表层指标}与\textbf{深层指标}两个层级：

\textbf{表层指标——回答"系统运行得如何"：}
\begin{enumerate}[label=(\arabic*)]
    \item \textbf{累计奖励（Cumulative Reward, CR）$\uparrow$：}每episode内全局奖励$R_t$的总和。
    \item \textbf{总冲突次数（Total Conflicts, TC）$\downarrow$：}每episode内冲突事件的总次数。
    \item \textbf{平均延迟（Average Delay, AD）$\downarrow$：}飞行器因速度偏离目标速度产生的运行延迟均值（s）。
\end{enumerate}

\textbf{深层指标——回答"系统做得有多好"：}
\begin{enumerate}[label=(\arabic*), resume]
    \item \textbf{通行完成率（Passing Completion Rate, PCR）$\uparrow$：}成功到达目标区域的飞行器数量占比。\textbf{注意：PCR仅统计到达结果，不区分到达过程中是否发生过冲突。}
    \item \textbf{无冲突完成率（Conflict-Free Completion Rate, CFCR）$\uparrow$：}全程零冲突且所有$N$架飞行器均完成通行的episode占比。\textbf{CFCR是PCR的子集。}
    \item \textbf{每架次平均冲突次数（Mean Conflicts per Aircraft per Episode, MCAE）$\downarrow$：}TC除以$N$，消除密度影响。
    \item \textbf{最小间隔违反率（Minimum Separation Violation Rate, MSVR）$\downarrow$：}存在安全距离违反的时间步数占比（\%）。
\end{enumerate}

\textbf{PCR与CFCR的本质区别：}PCR $= 99.1\%$与CFCR $= 68.4\%$之间30.7个百分点的差值意味着约31\%的episode中，所有飞行器虽最终到达了目标区域，但途中至少经历了一次安全距离违反——飞行器在有冲突的情况下通过调整航迹"挣扎着通过了"。若仅报告PCR（99.1\%），读者将获得系统"几乎完美"的错误印象；CFCR的引入揭示了这一表象下的真实安全状态。

\subsection{训练收敛性分析}

模型训练1000轮的收敛过程可划分为三个特征阶段：
（I）快速收敛期（第0--200轮）——累计奖励从约$-2000$上升至约$+800$；
（II）渐进优化期（第200--600轮）——奖励波动上升至约$+1100\sim1300$，课程学习触发多次难度进阶；
（III）精细收敛期（第600--1000轮）——奖励稳定在$+1286.7 \pm 18.2$，冲突次数收敛至$18.2 \pm 2.1$。

\begin{figure}[H]
    \centering
    \includegraphics[width=0.8\textwidth]{QQ图片20260504231524.jpg}
    \caption{训练收敛曲线——包含累计奖励（左上）、冲突次数（右上）、平均延迟（左下）与PCR/CFCR（右下）随训练轮次（0--1000）的变化。阴影区域表示5次独立训练的$\pm1$标准差。}
    \label{fig:training_curve}
\end{figure}

\subsection{综合性能对比}

$N=21$架eVTOL测试场景下的综合性能对比如表\ref{tab:comparison}所示。所有方法均报告5次独立测试的均值$\pm$标准差。

\begin{table}[H]
    \centering
    \caption{不同方法的综合性能对比（21架eVTOL，5次独立测试，均值$\pm$标准差）}
    \label{tab:comparison}
    \begin{tabular}{|c|c|c|c|c|}
        \hline
        \textbf{方法（类别）} & \textbf{CR$\uparrow$} & \textbf{TC$\downarrow$} & \textbf{AD(s)$\downarrow$} & \textbf{PCR(\%)$\uparrow$} \\
        \hline
        RR（传统规则） & $-1256.3 \pm 89.2$ & $42.6 \pm 5.3$ & $1.58 \pm 0.12$ & $89.2 \pm 3.5$ \\
        FCFS（传统规则） & $-1187.5 \pm 76.4$ & $38.9 \pm 4.7$ & $1.42 \pm 0.10$ & $91.5 \pm 2.8$ \\
        VO（工程安全） & $-512.3 \pm 45.6$ & $27.5 \pm 3.5$ & $0.92 \pm 0.08$ & $94.2 \pm 2.0$ \\
        ORCA（工程安全） & $-198.7 \pm 35.8$ & $22.8 \pm 2.9$ & $0.76 \pm 0.06$ & $96.0 \pm 1.6$ \\
        集中式DQN（单RL） & $-892.5 \pm 65.3$ & $28.3 \pm 3.9$ & $1.21 \pm 0.09$ & $94.7 \pm 2.1$ \\
        独立A2C（独立ML） & $-657.8 \pm 52.6$ & $19.7 \pm 2.8$ & $0.93 \pm 0.07$ & $96.3 \pm 1.7$ \\
        MADDPG（CTDE） & $-215.6 \pm 41.2$ & $25.4 \pm 3.2$ & $0.87 \pm 0.06$ & $95.8 \pm 1.9$ \\
        COMA（CTDE） & $187.3 \pm 35.7$ & $22.1 \pm 2.9$ & $0.79 \pm 0.06$ & $97.1 \pm 1.5$ \\
        QMIX（CTDE） & $426.5 \pm 28.9$ & $23.1 \pm 2.7$ & $0.76 \pm 0.05$ & $97.5 \pm 1.3$ \\
        MAPPO（CTDE） & $892.4 \pm 23.5$ & $20.3 \pm 2.4$ & $0.65 \pm 0.04$ & $98.2 \pm 1.1$ \\
        \textbf{本文方法} & $\mathbf{1286.7 \pm 18.2}$ & $\mathbf{18.2 \pm 2.1}$ & $\mathbf{0.57 \pm 0.03}$ & $\mathbf{99.1 \pm 0.8}$ \\
        \hline
    \end{tabular}
\end{table}

从表\ref{tab:comparison}可以得出以下层次化的分析：

\textbf{第一层：学习方法相较于规则方法的压倒性优势。}本文方法相较RR规则基线：冲突次数降低57.2\%（42.6$\to$18.2），平均运行延迟缩短63.9\%（1.58$\to$0.57 s），PCR提升9.9个百分点（89.2\%$\to$99.1\%）。规则方法虽有计算成本极低（$<1$ ms/步）的优势，但其性能上限受限于固定调度逻辑的"盲目性"——无法感知飞行器间的具体相对运动状态。

\textbf{第二层：工程安全方法（VO/ORCA）的性能定位。}ORCA在TC（22.8次）和PCR（96.0\%）上显著优于传统规则方法（RR: 42.6次，FCFS: 38.9次），验证了分布式几何避障在低空交通场景中的有效性。然而，ORCA的性能仍显著低于本文方法（TC: 18.2 vs.\ 22.8, CFCR: 68.4\% vs.\ 约50\%），主要原因在于：（a）ORCA仅基于当前几何关系决策，缺乏对未来轨迹的预测性建模；（b）在高密度交互场景中，速度空间的可行解可能退化，导致频繁的"震荡"行为。ORCA与MARL方法的混合架构——ORCA负责战术级紧急避障、MARL负责战略级协同优化——是一个有前景的研究方向。

\textbf{第三层：CTDE-MARL方法内部的架构差异。}QMIX受限于单调性约束，在$N$动态变化场景中性能受限。COMA的反事实基线在高维连续动作空间中估计方差较大。MADDPG受限于确定性策略梯度在高维交互场景中的探索不足。MAPPO是性能最强的CTDE-MARL基线。

\textbf{第四层：与MAPPO的差距——数值温和但结构重要。}本文方法相较MAPPO的改进在数值上是温和的：冲突次数降低10.3\%（20.3$\to$18.2），延迟缩短12.3\%（0.65$\to$0.57 s）。这一增量主要来自DyGAT的三个组件协同——空间注意力、时序注意力与物理量驱动动态边权重。

\subsubsection{安全指标专项对比}

表\ref{tab:safety_metrics}报告了四项安全攸关指标的对比结果。

\begin{table}[H]
    \centering
    \caption{安全指标专项对比（21架eVTOL，5次独立测试，均值$\pm$标准差）}
    \label{tab:safety_metrics}
    \begin{tabular}{|c|c|c|c|c|}
        \hline
        \textbf{方法} & \textbf{CFCR(\%)$\uparrow$} & \textbf{MCAE$\downarrow$} & \textbf{MSVR(\%)$\downarrow$} & \textbf{PCR(\%)$\uparrow$} \\
        \hline
        RR & $12.3 \pm 4.1$ & $2.03 \pm 0.25$ & $8.7 \pm 1.2$ & $89.2 \pm 3.5$ \\
        FCFS & $18.7 \pm 3.6$ & $1.85 \pm 0.22$ & $7.4 \pm 1.0$ & $91.5 \pm 2.8$ \\
        VO & $42.0 \pm 3.5$ & $1.31 \pm 0.17$ & $4.8 \pm 0.7$ & $94.2 \pm 2.0$ \\
        ORCA & $50.2 \pm 3.0$ & $1.09 \pm 0.14$ & $3.8 \pm 0.6$ & $96.0 \pm 1.6$ \\
        集中式DQN & $37.5 \pm 3.2$ & $1.35 \pm 0.19$ & $5.1 \pm 0.8$ & $94.7 \pm 2.1$ \\
        独立A2C & $48.2 \pm 2.8$ & $0.94 \pm 0.13$ & $3.6 \pm 0.6$ & $96.3 \pm 1.7$ \\
        MADDPG & $42.6 \pm 3.0$ & $1.21 \pm 0.15$ & $4.2 \pm 0.7$ & $95.8 \pm 1.9$ \\
        COMA & $51.3 \pm 2.7$ & $1.05 \pm 0.14$ & $3.8 \pm 0.6$ & $97.1 \pm 1.5$ \\
        QMIX & $53.7 \pm 2.5$ & $1.10 \pm 0.13$ & $3.5 \pm 0.5$ & $97.5 \pm 1.3$ \\
        MAPPO & $57.8 \pm 2.2$ & $0.97 \pm 0.11$ & $3.0 \pm 0.4$ & $98.2 \pm 1.1$ \\
        \textbf{本文方法} & $\mathbf{68.4 \pm 1.8}$ & $\mathbf{0.87 \pm 0.10}$ & $\mathbf{2.4 \pm 0.3}$ & $\mathbf{99.1 \pm 0.8}$ \\
        \hline
    \end{tabular}
\end{table}

表\ref{tab:safety_metrics}揭示了PCR与CFCR之间的巨大鸿沟——本文方法PCR $= 99.1\%$但CFCR仅$68.4\%$，两者相差30.7个百分点。ORCA的CFCR约为50.2\%，虽优于规则方法（RR: 12.3\%），但显著低于学习方法——这是因为ORCA的纯几何避障策略缺乏对未来多步轨迹的规划能力，在21架飞行器的高密度场景中经常导致"先紧急避障、后恢复航线"的震荡行为。MAPPO的PCR $= 98.2\%$与CFCR $= 57.8\%$之间也存在40.4个百分点的差距，表明这一问题是CTDE-MARL方法的共性挑战。

\subsubsection{冲突严重程度分级分析}

CFCR $= 68.4\%$意味着约31\%的episode仍有冲突。为回答"剩余31\%的冲突有多严重"，本文将安全距离违反按$d_{\min}/d_{\text{safety}}$比值分为三个严重程度等级：

\begin{itemize}
    \item \textbf{轻微违规（$0.8 \leq d_{\min}/d_{\text{safety}} < 1.0$）：}瞬时、小幅度的安全距离偏离，通常由传感器噪声或风扰导致的短时超调引起。
    \item \textbf{中等违规（$0.5 \leq d_{\min}/d_{\text{safety}} < 0.8$）：}显著的安全距离违反，表明飞行器间存在真实的碰撞风险。
    \item \textbf{严重违规/NMAC（$d_{\min}/d_{\text{safety}} < 0.5$）：}接近空中碰撞（Near Mid-Air Collision），在任何运营标准下均不可接受。ICAO将NMAC定义为飞行器间距离小于500 ft（约152 m），在本文场景中对应$d_{\min}/d_{\text{safety}} < 0.5$（$d_{\text{safety}} \approx 85$ m时对应$d_{\min} < 42.5$ m）。
\end{itemize}

表\ref{tab:severity}报告了50次测试episode中所有安全距离违反事件的严重程度分布。

\begin{table}[H]
    \centering
    \caption{冲突严重程度分级统计（50次测试episode，21架eVTOL）}
    \label{tab:severity}
    \begin{tabular}{|c|c|c|c|}
        \hline
        \textbf{$d_{\min}/d_{\text{safety}}$范围} & \textbf{事件数} & \textbf{占比} & \textbf{严重程度} \\
        \hline
        $[0.8, 1.0)$ & $\sim1250$ & $\sim80\%$ & 轻微入侵 \\
        $[0.5, 0.8)$ & $\sim250$ & $\sim16\%$ & 中等接近 \\
        $<0.5$（NMAC级别） & $\sim15$ & $<1\%$ & 严重接近 \\
        \hline
    \end{tabular}
\end{table}

在31\%的存在冲突的episode中，约80\%的冲突属于安全距离轻微入侵（$d_{\min}/d_{\text{safety}} \in [0.8, 1.0)$），NMAC级别（$d_{\min} < 0.5d_{\text{safety}} \approx 42.5$ m）的严重接近事件占比不足1\%。将CFCR从68.4\%提升至95\%以上、同时消除NMAC级别事件是下一步的核心优化目标。

\subsection{统计显著性检验}

对本文方法与MAPPO在四项关键指标上的性能差异进行Welch双样本$t$检验（双侧，不等方差，$n_1 = n_2 = 5$），结果如表\ref{tab:statistical_tests}所示。

\begin{table}[H]
    \centering
    \caption{本文方法 vs.\ MAPPO的统计显著性检验（$n=5$，双侧Welch $t$检验，$\alpha=0.05$）}
    \label{tab:statistical_tests}
    \begin{tabular}{|c|c|c|c|c|c|}
        \hline
        \textbf{指标} & \textbf{均值差} & \textbf{$t$统计量} & \textbf{$p$值} & \textbf{Cohen's $d$} & \textbf{显著性} \\
        \hline
        TC & $2.1 \downarrow$ & $t(8)=1.47$ & $p=0.180$ & $0.93$ & 不显著（n.s.） \\
        CFCR & $10.6\% \uparrow$ & $t(8)=8.34$ & $p<0.001$ & $5.27$ & \textbf{高度显著} \\
        MCAE & $0.10 \downarrow$ & $t(8)=1.50$ & $p=0.172$ & $0.95$ & 不显著（n.s.） \\
        MSVR & $0.6\% \downarrow$ & $t(8)=2.68$ & $p=0.029$ & $1.69$ & \textbf{显著} \\
        \hline
    \end{tabular}
\end{table}

\subsubsection{统计效力的局限性说明}

本文的统计检验基于$n=5$次独立运行。对于TC指标，尽管Cohen's $d \approx 0.93$属于大效应量（$>0.8$），但由于样本量限制，Welch $t$检验的统计效力（power）仅约38\%——这意味着即使存在真实差异，当前样本量也只有38\%的概率检测到统计显著性。要达到80\%的统计效力（$\alpha=0.05$，$d=0.93$），每组需要至少$n \geq 20$次独立运行。因此，本文报告的TC"无显著差异"（$p=0.180$）应理解为"在$n=5$条件下未检测到显著差异"，而非"两组真实均值相等"的证据。

相比之下，CFCR的效应量极大（$d \approx 5.27$），即使$n=5$也达到了$>99\%$的统计效力，结论高度可信。MSVR的效应量$d \approx 1.69$（大效应量），统计效力约70\%——在$\alpha=0.05$水平上的显著性（$p=0.029$）具有中等可信度。

在更大规模的计算资源下重复实验以提升TC检验的统计效力，是下一步工作之一。

\subsection{消融实验}

表\ref{tab:ablation}报告了各核心模块的独立性能贡献。消融实验遵循"一次移除一个模块，其余保持不变"的单变量控制原则。

\begin{table}[H]
    \centering
    \caption{消融实验结果（21架eVTOL，5次独立测试，均值$\pm$标准差）}
    \label{tab:ablation}
    \resizebox{0.92\textwidth}{!}{%
    \begin{tabular}{|c|c|c|c|c|}
        \hline
        \textbf{模型配置} & \textbf{CR$\uparrow$} & \textbf{TC$\downarrow$} & \textbf{AD(s)$\downarrow$} & \textbf{PCR(\%)$\uparrow$} \\
        \hline
        完整模型（本文） & $\mathbf{1286.7 \pm 18.2}$ & $\mathbf{18.2 \pm 2.1}$ & $\mathbf{0.57 \pm 0.03}$ & $\mathbf{99.1 \pm 0.8}$ \\
        \hline
        $-$ DyGAT（改用标准GAT） & $215.3 \pm 32.6$ & $49.8 \pm 4.5$ & $0.79 \pm 0.05$ & $94.3 \pm 1.9$ \\
        \hline
        $-$ 课程学习（直接$N=21$训练） & $587.2 \pm 27.4$ & $28.4 \pm 3.1$ & $0.68 \pm 0.04$ & $97.2 \pm 1.4$ \\
        \hline
        $-$ 冲突轨迹加权（$\lambda \equiv 1.0$） & $763.5 \pm 24.1$ & $24.6 \pm 2.8$ & $0.65 \pm 0.04$ & $97.8 \pm 1.2$ \\
        \hline
        $-$ 课程学习 $-$ 冲突轨迹加权 & $128.6 \pm 38.5$ & $35.7 \pm 3.9$ & $0.82 \pm 0.06$ & $95.6 \pm 1.7$ \\
        \hline
    \end{tabular}%
    }
\end{table}

\textbf{DyGAT是性能贡献最大的单一模块。}移除DyGAT（替换为标准GAT）后，冲突次数从18.2激增至49.8，相对增幅达173.6\%，确定性地表明DyGAT的时空耦合建模是本文方法性能的\textbf{绝对基石}。

\textbf{课程学习与冲突轨迹加权形成互补协同。}单独移除课程学习导致冲突上升56.0\%，单独移除轨迹加权导致冲突上升35.2\%。两者同时移除时冲突上升96.2\%，该增幅与单独效应之和基本一致，表明两种机制的作用路径大部分相互独立。

\subsection{密度鲁棒性与敏感性分析}

\subsubsection{飞行器密度鲁棒性}

本文方法在$N \in \{5, 10, 15, 21\}$四种密度场景下的冲突次数增长斜率（$\approx 0.75$次/架）显著低于MAPPO（$\approx 1.15$次/架）与RR（$\approx 2.4$次/架）。DyGAT在密度升高时优势愈发显著的核心机制在于：飞行器越多，精准区分不同飞行器对的冲突风险等级越关键——将有限的注意力资源集中于高风险交互对。

\begin{figure}[H]
    \centering
    \includegraphics[width=0.65\textwidth]{QQ图片20260524171218.png}
    \caption{不同飞行器密度（$N=5, 10, 15, 21$）下的冲突次数对比。本文方法、MAPPO与RR三类方法分别展示。误差条表示$\pm1$标准差（5次独立测试）。}
    \label{fig:density_robustness}
\end{figure}

\subsubsection{超参数敏感性分析}

\textbf{基础安全距离 $d_{\text{base}}$（40--70 m）：}$d_{\text{base}}$从40 m增至70 m时，冲突次数从22.5降至14.7（$\downarrow 34.7\%$），但延迟从0.52 s升至0.63 s（$\uparrow 21.2\%$），体现了航空交通管理中经典的安全-效率基本权衡。

\textbf{风扰强度（0--0.3 m/s）：}风扰从0增至0.3 m/s时，冲突次数上升约27.6\%。性能退化幅度温和——DyGAT的时序注意力通过学习历史轨迹的时序趋势，部分补偿了风扰带来的位置预测不确定性。但需指出：当前均匀随机风扰模型是对真实低空湍流的显著简化，真实Dryden风场下因空间相干结构导致的性能退化可能更大。

\textbf{传感器噪声 $\sigma$（0--1.5 m）：}$\sigma$从0增至1.5 m时，冲突次数仅上升约14.6\%——所有参数中影响最小。LSTM编码器对连续10步历史状态序列的隐式平滑作用有效过滤了高频观测噪声。

\begin{figure}[H]
    \centering
    \includegraphics[width=0.8\textwidth]{QQ图片20260524162300.png}
    \caption{超参数敏感性分析——分别改变基础安全距离$d_{\text{base}}$（左）、风扰强度（中）与传感器噪声$\sigma$（右）时冲突次数与平均延迟的变化。每个数据点为5次独立测试的均值，误差条为$\pm1$标准差。红色虚线标记默认参数值。}
    \label{fig:sensitivity_analysis}
\end{figure}

\subsection{计算成本分析}

本文方法在$N=21$时平均推理时间为$4.5$ ms（NVIDIA RTX 4060 GPU），远低于仿真步长$500$ ms——实时性裕度约为$111\times$。计算时间从$N=5$时的$1.2$ ms增长至$N=21$时的$4.5$ ms，增长斜率约为$0.21$ ms/架，在所有CTDE-MARL方法中最低。

\begin{figure}[H]
    \centering
    \includegraphics[width=0.8\textwidth]{QQ图片20260524163903.jpg}
    \caption{不同飞行器数量下的平均单步推理计算时间对比（NVIDIA RTX 4060 GPU）。本文方法、MAPPO、QMIX与COMA四类CTDE-MARL方法展示。}
    \label{fig:computation_cost}
\end{figure}

%==========================================================================
% 5. 讨论
%==========================================================================
\section{讨论}

\subsection{核心贡献的机制解读}

本文方法的性能优势可归因于三个设计要素的协同：（1）DyGAT的精准时空耦合刻画是性能基石（消融中移除导致冲突暴增173.6\%）；（2）课程学习与冲突轨迹加权形成训练机制协同（两模块作用路径高度独立）；（3）GAT-Mixer的架构适配性通过改善训练信号质量间接提升策略质量。

\subsection{实际应用价值}

尽管处于proof-of-concept阶段，本文模型在以下方面展现了应用潜力：

\textbf{（1）模块化可集成性。}模型的分布式执行架构意味着每架eVTOL仅需本地观测即可做出决策，4.5 ms的推理延迟（RTX 4060）远低于UTM系统典型的50--200 ms决策周期，满足实时性要求。

\textbf{（2）训练-部署分离。}集中训练可利用地面高性能计算集群完成，训练好的策略网络（约2--5 MB参数量）可直接部署到机载边缘计算设备，不依赖训练阶段的全局状态和通信基础设施。

\textbf{（3）环境适应性。}模型在0--0.3 m/s风扰和0--1.5 m传感器噪声范围内表现出稳定的性能衰减特性（冲突增加$<28\%$），说明其对真实低空环境的不确定性具有一定容忍度。

上述应用潜力的兑现，依赖于第5.4节所述局限性的逐步解决。

\subsection{与MS-GRL的详细比较}

Li等提出的MS-GRL\cite{ref24}与本文目标最为接近。两者在技术路线上的本质差异决定了互补而非替代的关系：

\textbf{（1）算法选择（Double DQN vs.\ PPO）。}MS-GRL采用Double DQN处理离散动作空间（速度档位选择）。离散动作的优势在于训练稳定，代价是控制粒度粗糙。本文采用PPO处理连续加速度控制，允许飞行器在物理约束范围内任意调节加速度，这在eVTOL的精细化速度调控需求中是更自然的选择。

\textbf{（2）图结构（固定阈值 vs.\ 物理量驱动动态权重）。}MS-GRL的多尺度GAT基于固定距离阈值。本文的DyGAT通过距离、速度差和航向差三个物理量计算连续动态边权重$w_{ij} \in (0, 1]$——两对距离同为200 m的飞行器，若一对同向等速飞行（$w_{ij} \approx 1.0$），另一对相向高速靠近（$w_{ij} \approx 0.015$），后者的注意力权重仅为前者的约1/67，风险区分能力远超固定阈值方案。

\textbf{（3）场景规模与针对性。}MS-GRL面向通用UAV编队场景（100架规模），强调大规模可扩展性。本文聚焦eVTOL交叉路口的特定运营场景（21架规模），在安全约束、物理建模与运营逻辑上更具针对性。

两者在图增强MARL这一技术路线上形成了互补——MS-GRL为大规模UAV网络的图结构设计提供了参考，本文为eVTOL精细化管控的物理量驱动动态图提供了方案。

\subsection{方法局限性}

本文方法存在以下主要局限性，按严重程度排序：

\textbf{（1）场景维度限制（严重程度：高）。}当前仿真场景为二维单交叉路口的简化模型。真实低空交通涉及三维高度层（含垂直起降）、多路口网络拓扑、起降场容量约束与通航走廊的空间约束。二维简化忽略了垂直维度的冲突规避——在真实三维空间中，飞行器可通过高度层分离规避冲突，而无需仅依赖水平面的速度/航向调整。

\textbf{（2）通信假设的理想化（严重程度：高）。}CTDE训练假设全局状态可无延迟、无丢失地获取。在真实UTM环境中，通信链路面临延迟（50--200 ms）、数据包丢失与有限带宽的约束。

\textbf{（3）飞行器同质化假设（严重程度：中）。}所有$N$架飞行器使用相同的运动学参数。实际AAM生态中，不同制造商的eVTOL具有差异化的飞行包线。

\textbf{（4）风扰与噪声建模的简化（严重程度：中）。}风扰采用均匀随机分布，噪声采用高斯白噪声。真实低空风场具有显著的空间相关性与时间连续性——Dryden模型（MIL-F-8785C）\cite{ref20}的湍流积分尺度可达数百米。

\textbf{（5）基线比较范围不足（严重程度：中）。}MPC和CBF方法尚未纳入实验对比。与这些工程安全关键方法的比较是投稿前或修改阶段的优先补充实验。

\textbf{（6）统计样本量限制（严重程度：中低）。}$n=5$的小样本量限制了对中等效应量差异的统计检测能力。需在更大规模的计算资源下扩大样本量至$n \geq 20$。

%==========================================================================
% 6. 结论与未来工作
%==========================================================================
\section{结论与未来工作}

本文针对城市低空十字交叉路口的eVTOL协同管控与冲突避免问题，提出了一种基于耦合多智能体强化学习的管控模型。模型核心设计——DyGAT动态图注意力网络、课程学习与冲突轨迹加权训练框架——分别在时空耦合建模与稀疏奖励处理两个关键维度上做出了针对性改进。

在21架eVTOL交叉路口仿真场景中，本文方法相较规则基线（RR）冲突次数降低57.2\%，平均延迟缩短63.9\%；相较当前最优CTDE-MARL方法MAPPO，冲突次数降低10.3\%（$p=0.180$，未达统计显著），但CFCR从57.8\%提升至68.4\%（$p<0.001$），MSVR从3.0\%降至2.4\%（$p=0.029$）。消融实验表明DyGAT是性能贡献最大的单一模块（移除后冲突增加173.6\%）。

本文的核心方法论贡献在于：（1）引入CFCR等深层安全指标，揭示了PCR（99.1\%）与CFCR（68.4\%）之间30.7个百分点的鸿沟；（2）通过Welch $t$检验与事后统计功效分析，诚实呈现了差异化统计图景——这种区分"统计功效不足"与"效应量为零"的分析范式，对于安全攸关的AI系统评估具有方法论示范价值。

本文的局限性（二维简化场景、均匀风扰模型、有限基线覆盖、$n=5$小样本量、冲突严重程度分布统计初步）在第5.4节中已有坦诚讨论。这些局限意味着当前结果应被解读为概念验证（proof-of-concept），而非生产就绪（production-ready）系统的性能声明。

\textbf{未来工作按优先级分为以下方向：}

\textbf{近期（6--12个月，投稿前优先补充）：}
\begin{enumerate}[label=(\arabic*)]
    \item 将$d_{\min}/d_{\text{safety}}$分级分析系统化——扩大测试样本量至50--100次episode，输出完整的三级分布直方图和每级占比的置信区间；
    \item 扩大独立运行次数至$n \geq 20$，提高TC等指标的统计检验效力至80\%以上；
    \item 完成DyGAT组件级消融——独立移除空间注意力、时序注意力与动态边权重，量化各子模块的独立贡献；
    \item 完成$\lambda \in \{1.0, 1.5, 2.0, 2.5, 3.0\}$的参数敏感性全网格分析，确定最优区间；
    \item 引入MPC/CBF作为基线或探索"MARL$+$CBF"混合架构——MARL负责战略级协同优化、CBF负责硬安全约束的实时执行。
\end{enumerate}

\textbf{中期（1--2年）：}
\begin{enumerate}[label=(\arabic*), resume]
    \item 将场景从二维单交叉路口扩展至三维多路口网络，引入垂直起降动力学与起降场容量约束；
    \item 引入通信延迟（50--200 ms）、数据包丢失（5--15\%丢包率）等实际网络约束，评估分布式执行架构的鲁棒性；
    \item 集成Dryden MIL-F-8785C的完整von K\'{a}rm\'{a}n风场模型（$L_u, L_v, L_w$三参数），替代当前的简化均匀风扰；
    \item 探索"MARL$+$ORCA"混合架构——MARL负责高层协同优化、ORCA负责战术级紧急避障，发挥两种方法的互补优势。
\end{enumerate}

\textbf{远期（2--3年）：}
\begin{enumerate}[label=(\arabic*), resume]
    \item 支持不同类型eVTOL（Joby S4、Volocopter VoloCity、Lilium Jet等）的异构动力学参数，研究异构群体中的策略泛化与自适应；
    \item 通过PX4/ArduPilot硬件在环仿真（Hardware-in-the-Loop, HIL）进行实飞验证——先以3--5架小型无人机在受控环境中验证策略的物理可执行性，再逐步扩展至更大规模和更复杂环境。
\end{enumerate}

代码将在论文接收后开源，以促进低空MARL管控领域的可复现研究。

%==========================================================================
% 参考文献
%==========================================================================
\begin{thebibliography}{99}

\bibitem{ref26} Federal Aviation Administration. Urban air mobility (UAM): concept of operations v2.0[R]. Washington, DC: FAA, 2023.

\bibitem{ref31} MENON P K, SWERIDUK G D, BILIMORIA K D. New approach to modeling, analysis, and control of air traffic flow[J]. Journal of Guidance, Control, and Dynamics, 2004, 27(5): 737-744.

\bibitem{ref27} RICHARDS A, HOW J P. Aircraft trajectory planning with collision avoidance using mixed integer linear programming[C]//Proceedings of the 2002 American Control Conference (ACC). Piscataway: IEEE, 2002: 1936-1941.

\bibitem{ref22} VAN DEN BERG J, GUY S J, LIN M, et al. Reciprocal n-body collision avoidance[C]//Robotics Research: The 14th International Symposium ISRR (Springer Tracts in Advanced Robotics, Vol.\ 70). Berlin: Springer, 2011: 3-19.

\bibitem{ref23} AMES A D, COOGAN S, EGERSTEDT M, et al. Control barrier functions: theory and applications[C]//2019 18th European Control Conference (ECC). Piscataway: IEEE, 2019: 3420-3431.

\bibitem{ref1} AGOGINO A, TUMER K. Regulating air traffic flow with coupled agents[C]//Proceedings of the 7th International Conference on Autonomous Agents and Multiagent Systems (AAMAS 2008). Richland: IFAAMAS, 2008: 535-542.

\bibitem{ref2} AGOGINO A K, TUMER K. A multiagent approach to managing air traffic flow[J]. Autonomous Agents and Multi-Agent Systems, 2012, 24(1): 1-25.

\bibitem{ref3} BRITTAIN M, WEI P. Autonomous separation assurance in a high-density en route sector: a deep multi-agent reinforcement learning approach[C]//2019 IEEE Intelligent Transportation Systems Conference (ITSC). Piscataway: IEEE, 2019: 3256-3262.

\bibitem{ref9} GHOSH S, LAGUNA S, LIM S H, et al. A deep ensemble method for multi-agent reinforcement learning: a case study on air traffic control[C]//Proceedings of the 31st International Conference on Automated Planning and Scheduling (ICAPS 2021). Palo Alto: AAAI Press, 2021: 468-476.

\bibitem{ref6} XIE Y B, GARDI A, SABATINI R. Reinforcement learning-based flow management techniques for urban air mobility and dense low-altitude air traffic operations[C]//2021 IEEE/AIAA 40th Digital Avionics Systems Conference (DASC). Piscataway: IEEE, 2021: 1-10.

\bibitem{ref24} LI X, WANG Y, ZHANG H, et al. Multiscale-graph-enhanced reinforcement learning for conflict resolution in dense UAV networks[J]. IEEE Internet of Things Journal, 2025, 12(22): 72134-72148.

\bibitem{ref7} SPATHARIS C, BASTAS A, KRAVARIS T, et al. Hierarchical multiagent reinforcement learning schemes for air traffic management[J]. Neural Computing and Applications, 2021, 33(14): 8649-8671.

\bibitem{ref10} ASLANI M, MESGARI M S, WIERING M. Adaptive traffic signal control with actor-critic methods in a real-world traffic network with different traffic disruption events[J]. Transportation Research Part C: Emerging Technologies, 2017, 85: 732-752.

\bibitem{ref11} 翟子洋, 郝茹茹, 董世浩. 大规模智慧交通信号控制中的强化学习和深度强化学习方法综述[J]. 计算机应用研究, 2024, 41(6): 1618-1627.

\bibitem{ref20} U.S. Department of Defense. Flying qualities of piloted airplanes: MIL-F-8785C[S]. Washington, DC: U.S. Department of Defense, 1980.

\bibitem{ref4} TUMER K, AGOGINO A. Distributed agent-based air traffic flow management[C]//Proceedings of the 6th International Joint Conference on Autonomous Agents and Multiagent Systems (AAMAS 2007). Honolulu: ACM, 2007: 330-337.

\bibitem{ref13} 王迎. 城市交通信号控制深度强化学习算法研究[D]. 济南: 山东交通学院, 2024.

\bibitem{ref16} 赖建辉. 基于深度强化学习的交通控制优化方法研究及实现[D]. 杭州: 浙江工业大学, 2020.

\bibitem{ref17} 张新武. 基于深度强化学习算法的交通信号控制优化研究[D]. 太原: 太原科技大学, 2024.

\bibitem{ref14} 赵晓华, 石建军, 李振龙, 赵国勇. 基于Q-learning和BP神经元网络的交叉口信号灯控制[J]. 公路交通科技, 2007, 24(7): 99-102.

\bibitem{ref15} 承向军, 常歆识, 杨肇夏. 基于Q-学习的交通信号控制方法[J]. 系统工程理论与实践, 2006(8): 136-140.

\bibitem{ref18} 卢守峰, 张术, 刘喜敏. 平均排队长度差最小的单交叉口在线Q学习模型[J]. 公路交通科技, 2014, 31(11): 116-122.

\bibitem{ref19} 江波, 李彦冬, 刘威陇, 等. 应用深度Q学习的机场道口协同智能调速方法[J]. 航空计算技术, 2022, 52(1): 26-30.

\bibitem{ref8} 张春阳, 席雅琪. 人工智能在城市交通控制中的应用综述[J]. 信息系统工程, 2024(07): 33-36.

\bibitem{ref12} GHAVAMZADEH M, MAHADEVAN S, MAKAR R. Hierarchical multi-agent reinforcement learning[J]. Autonomous Agents and Multi-Agent Systems, 2006, 13(2): 197-229.

\bibitem{ref5} TUMER K, AGOGINO A. Improving air traffic management with a learning multiagent system[J]. IEEE Intelligent Systems, 2009, 24(1): 18-21.

\bibitem{ref28} Joby Aviation. Joby S4 technical specifications[EB/OL]. 2024. \url{https://www.jobyaviation.com/}

\bibitem{ref29} Volocopter GmbH. VoloCity technical specifications[EB/OL]. 2024. \url{https://www.volocopter.com/}

\bibitem{ref30} Lilium N.V. Lilium Jet technical specifications[EB/OL]. 2024. \url{https://www.lilium.com/}

\bibitem{ref25} CONG W, ZHANG S, CHEN J, et al. Do we really need complicated model architectures for temporal networks?[C]//Proceedings of the 11th International Conference on Learning Representations (ICLR 2023). Kigali: ICLR, 2023.

\bibitem{ref21} LIANG X, MA Y, FENG Y, et al. PTR-PPO: proximal policy optimization with prioritized trajectory replay[EB/OL]. arXiv preprint arXiv:2112.03798, 2021.

\end{thebibliography}

\end{document}
